\documentclass[a4paper,11pt]{letter}
\usepackage[utf8]{inputenc}
\usepackage[T1]{fontenc}
\usepackage{lmodern}
\usepackage[english]{babel}
\usepackage[colorlinks=true, urlcolor=blue]{hyperref}
\usepackage{geometry}
\geometry{left=1in, right=1in, top=1in, bottom=1in}

\name{Baha Khammari}
\address{D\"usseldorf, Germany \\
    \href{mailto:b.e.khammari@gmail.com}{b.e.khammari@gmail.com} \\
    +49\,1520\,9016956}

\begin{document}

\begin{letter}{%
    D\"usseldorf Graduate School of Economics (DGSE) \\
    D\"usseldorf Institute for Competition Economics (DICE) \\
    Heinrich-Heine-Universit\"at D\"usseldorf \\
    Universit\"atsstra\ss e 1, 40225 D\"usseldorf
}

\opening{Dear Members of the Selection Committee,}

I am writing to apply for the doctoral programme at the D\"usseldorf Graduate School of Economics. I am completing my M.Sc.\ in Economics at the University of Cologne (expected September 2026, current GPA: 1.6/1.0) and will be available from October 2026. My research interests are in applied microeconomics, causal inference, and education/labour economics, with a methodological focus on difference-in-differences estimation and heterogeneous treatment effects.

My doctoral project asks which higher education access mechanism, whether scholarships, loans, or affirmative-action quotas, is most effective at reducing labour market inequality, and for whom. I use Brazil as an empirical laboratory because it implemented all three at national scale within a fifteen-year window, but the research questions and comparative framework are designed to inform the same policy choices in Germany, the EU, and other countries.

My master's thesis (grade: 1.0/1.0, supervised by Anna Person, M.Sc.) provides the first causal analysis of Brazil's ProUni scholarship programme on formal-sector wages at the microregion level. I constructed a panel of 558 regions over 2005--2019 from administrative microdata ($\sim$5\,million RAIS observations) and applied the Callaway \& Sant'Anna (2021) heterogeneity-robust staggered DiD estimator with doubly-robust inference. The key finding, a non-monotonic dose--response with significant wage gains in low-exposure regions that vanish at higher intensities, was entirely masked by conventional two-way fixed effects. The project required managing large-scale administrative datasets via SQL/BigQuery and implementing recent econometric methods in R.

Two members of the DICE faculty work on questions that intersect directly with my project. Prof.\ Dr.\ Hannah Schildberg-H\"orisch's research on the economics of education, in particular how socioeconomic status shapes skill formation and educational outcomes, addresses the upstream side of the inequality my project traces into the labour market. Her work on school-based interventions and children's economic preferences complements my focus on the downstream labour market effects of access policies. Combining both perspectives would strengthen the causal chain from education access to labour market outcomes.

Prof.\ Dr.\ Jens S\"udekum's research on the causal labour market effects of policy shocks at the regional level, including his DFG-funded project on trade integration, job stability, and geographical labour mobility, shares methodological common ground with my own approach: both use regional variation in policy exposure and administrative data to identify causal effects on worker outcomes. His expertise in regional economics would also inform my analysis of how policy effects vary across local labour markets.

Before the thesis, I worked for two years at PwC Germany (Assurance Solutions), where I built variance analyses on large financial datasets. An Erasmus+ semester at RSM Erasmus University (Rotterdam) added coursework in asset pricing and derivatives valuation.

My academic CV, a research proposal, and a thesis abstract are enclosed. References are available on request.

Thank you for considering my application.

\closing{Yours sincerely,}

\end{letter}

\end{document}
